\documentclass[11pt]{article}
% Shared preamble for Theseus user guides.
% Compiled with pdflatex. See docs/guides/BUILD.md for rationale.

\usepackage[a4paper,margin=0.85in]{geometry}
\usepackage[T1]{fontenc}
\usepackage[utf8]{inputenc}

% Body: Palatino (via mathpazo). Headings: Helvetica (sans).
% Palatino is requested in the house style as a substitute for Charter
% because Charter is not part of the standard TeX Live distribution
% that the build container ships.
\usepackage{mathpazo}
\usepackage[scaled=0.92]{helvet}
\usepackage{courier}
\renewcommand{\familydefault}{\rmdefault}

\usepackage{microtype}
\usepackage{parskip}
\usepackage{enumitem}
\usepackage{booktabs}
\usepackage{longtable}
\usepackage{titlesec}
\usepackage{xcolor}
\usepackage{fancyhdr}
\usepackage{fancyvrb}
\usepackage{graphicx}
% Allow `_` in text mode without escaping. The screenshot file names and
% some environment-variable references in the guides contain underscores.
\usepackage{underscore}
\usepackage{tcolorbox}
\tcbuselibrary{breakable,skins}
\usepackage[hidelinks,colorlinks=true,linkcolor=black,urlcolor=teal,citecolor=black]{hyperref}

% Sans for chapter and section headings; serif body keeps prose readable
% on screen and in print.
\titleformat{\section}{\sffamily\Large\bfseries}{\thesection}{0.7em}{}
\titleformat{\subsection}{\sffamily\large\bfseries}{\thesubsection}{0.6em}{}
\titleformat{\subsubsection}{\sffamily\normalsize\bfseries}{\thesubsubsection}{0.5em}{}

\pagestyle{fancy}
\fancyhf{}
\fancyhead[L]{\small\sffamily Theseus User Guide}
\fancyhead[R]{\small\sffamily\thepage}
\renewcommand{\headrulewidth}{0.3pt}

\setlength{\parindent}{0pt}

% Convenience commands shared by every guide.
\newcommand{\page}[1]{\textbf{#1}}
\newcommand{\file}[1]{\texttt{#1}}
\newcommand{\route}[1]{\texttt{#1}}
\newcommand{\env}[1]{\texttt{#1}}
\newcommand{\code}[1]{\texttt{#1}}

% Callout boxes used by every guide.
\newtcolorbox{tldr}{
  colback=teal!5!white,colframe=teal!55!black,
  fonttitle=\sffamily\bfseries,title=TL;DR,
  breakable,boxrule=0.4pt,arc=2pt
}
\newtcolorbox{youwillneed}{
  colback=gray!4!white,colframe=gray!55!black,
  fonttitle=\sffamily\bfseries,title=You will need,
  breakable,boxrule=0.4pt,arc=2pt
}
\newtcolorbox{notcovered}{
  colback=orange!5!white,colframe=orange!55!black,
  fonttitle=\sffamily\bfseries,title=This guide does NOT cover,
  breakable,boxrule=0.4pt,arc=2pt
}
\newtcolorbox{preview}{
  colback=yellow!8!white,colframe=yellow!55!black,
  fonttitle=\sffamily\bfseries,title=Preview --- partially shipped,
  breakable,boxrule=0.4pt,arc=2pt
}
\newtcolorbox{nextup}{
  colback=blue!4!white,colframe=blue!55!black,
  fonttitle=\sffamily\bfseries,title=Where to read next,
  breakable,boxrule=0.4pt,arc=2pt
}

% Insert a screenshot only when the file exists, so the build does not
% break before the Playwright capture run is wired up.
\newcommand{\screenshot}[3]{%
  \IfFileExists{screenshots/#1}{%
    \begin{figure}[h]\centering
      \includegraphics[width=\linewidth]{screenshots/#1}
      \caption{#2}\label{fig:#3}
    \end{figure}%
  }{%
    \begin{center}\small\itshape%
      [screenshot \texttt{screenshots/#1} not yet captured --- #2]%
    \end{center}%
  }%
}


\title{\sffamily\bfseries Currents, Articles, and Public Engagement\\[2pt]\large Guide 4 of 6}
\date{}
\author{}

\begin{document}
\maketitle
\thispagestyle{empty}

\begin{tldr}
Currents is the firm's live commentary surface. It pulls public events
in from X, decides whether the firm's corpus has anything substantive
to say, and either writes a citation-grounded opinion or abstains.
Every published opinion respects a set of structural invariants that
prevent it from speaking outside the corpus or revealing private
material. Beyond Currents proper, this guide explains the public-
article dispatcher (long-form essays built from clustered opinions),
the reader-response and critique-submission channels, the subscriber
digest, and how to read everything the firm publishes. Operator-only
actions (revocation, kill switches) are described in Guide 6.
\end{tldr}

\begin{youwillneed}
\begin{itemize}[leftmargin=*,itemsep=0.2em,topsep=0.2em]
  \item A signed-in Codex session, or an anonymous browser session
        for the public \route{/currents} page.
  \item Familiarity with how principles and conclusions work --- see
        Guide 2.
\end{itemize}
\end{youwillneed}

\begin{notcovered}
\begin{itemize}[leftmargin=*,itemsep=0.2em,topsep=0.2em]
  \item Operator-only overrides (revocation, kill switch, hold for
        review) --- Guide 6.
  \item The Oracle's question-answering surface --- Guide 3.
\end{itemize}
\end{notcovered}

\section{What Currents is}

Currents is a separate origin from the main Codex. The public page is
\route{/currents}; the founder-facing read view lives inside the Codex
under \route{/(authed)/founder-currents}. Each ``Current'' is one
\page{EventOpinion} row written by the firm in response to one outside
event.

The site keeps two kinds of public artifacts in this neighborhood:

\begin{itemize}[leftmargin=*,itemsep=0.2em]
  \item \page{Currents opinions} --- short, source-grounded
        statements attached to one specific event. They land on
        \route{/currents/<id>}.
  \item \page{Articles} --- long-form essays the firm publishes after
        clustering enough related opinions and material. They land on
        \route{/post/<slug>} (for blog-style articles tied to an
        upload) or \route{/c/<slug>/v/<version>} (for reviewed
        conclusions promoted to article shape).
\end{itemize}

\screenshot{04_currents_public_feed.png}{The public \route{/currents} feed, with two opinions visible and one event in abstention.}{cu-feed}

\section{The pipeline, top to bottom}

The Currents scheduler runs continuously, on a short cycle. Each
cycle:

\begin{enumerate}[leftmargin=*,itemsep=0.2em]
  \item \textbf{Discovery.} The system looks for recent high-engagement
        public posts on X.
  \item \textbf{Significance floor.} Each candidate post has a
        significance score computed from impressions, retweets,
        likes, replies, and quote-and-bookmark counts (a weighted
        log-sum). Posts below the configured floor are dropped.
  \item \textbf{Persistence and dedupe.} Survivors become event rows.
  \item \textbf{Enrichment.} The post is embedded, near-duplicate
        matched against recent events, and tagged with topic hints.
  \item \textbf{Relevance gate.} The event's fingerprint is compared
        against the firm's stored conclusions and claims. The gate
        requires at least two qualifying hits above a relevance
        threshold. If the corpus can't supply them, the event becomes
        an \emph{abstention} (the public page says ``the firm does
        not have enough material to comment'') and does not become an
        opinion.
  \item \textbf{Opinion generation.} The retrieval bundle is wrapped
        in unambiguous markers and handed to a language model with a
        strict contract: the model must cite each source it leans on
        by id and quote verbatim spans.
  \item \textbf{Citation validation.} Every quoted span in the model's
        output is checked against the cited source --- character by
        character. A fabrication anywhere in the output causes the
        whole opinion to abstain.
  \item \textbf{Strawman detection.} A separate guard checks that the
        opinion isn't mischaracterizing the source post.
  \item \textbf{Publication.} The opinion row is written; the public
        \route{/currents} page picks it up immediately.
  \item \textbf{Article dispatch.} A separate, longer-cycle job
        periodically scans recent opinions, clusters related ones,
        and --- if the cluster is dense enough --- drafts a long-form
        public article. The dispatcher is capped at a fixed number of
        articles per week.
\end{enumerate}

\section{The eight invariants}

A published Currents opinion must satisfy all eight of the following
structural invariants. They are enforced by the generator and the
publish path; an opinion that fails any one of them is rejected
before it reaches the public feed.

\begin{enumerate}[leftmargin=*,itemsep=0.3em]
  \item \textbf{Grounded.} Every non-trivial claim in the opinion
        text is backed by an inline \texttt{[C:<id>]} citation to a
        firm Conclusion or Principle. Free-floating assertions are
        not allowed.
  \item \textbf{Source-respecting.} Citations resolve only to
        publicly-safe corpus surfaces. Private transcripts and
        documents may inform internal reasoning but must not appear
        as public citation targets.
  \item \textbf{Quorum.} At least two relevance-cleared firm
        conclusions back the opinion. If the corpus cannot supply
        two, the system abstains rather than under-cite.
  \item \textbf{Voice-consistent.} The opinion is written in the
        firm's first-person voice rather than recapping or
        paraphrasing the source post.
  \item \textbf{Event-anchored.} Each opinion is bound to exactly one
        event. Multi-event commentary lives in the articles pipeline,
        not in Currents.
  \item \textbf{Revocable.} Every opinion is publishable iff
        \texttt{revokedAt IS NULL}. Operator revocation is the
        canonical inverse of publication; there is no separate
        ``hide'' flag.
  \item \textbf{Abstention-honest.} When the relevance gate fails,
        the system records a reason on the event and surfaces ``the
        firm does not yet have enough material to comment.''
        Abstentions are a public stance.
  \item \textbf{Provenance-complete.} Each opinion ships an audit
        trail: the source post, the relevance scores, the conclusions
        cited, and the generator parameters. The audit trail is
        addressable from the public page.
\end{enumerate}

\section{Reading the audit trail}

Open any opinion at \route{/currents/<id>}. Below the body you will
find:

\begin{itemize}[leftmargin=*,itemsep=0.2em]
  \item \textbf{Source post.} The original X post the opinion
        responds to. Always present, always external.
  \item \textbf{Cited conclusions.} A list of \texttt{[C:<id>]}
        references, each click-through to the firm's Conclusion on
        the public site. Private conclusions never appear here.
  \item \textbf{Relevance scores.} The score each cited conclusion
        had against the source event at generation time.
  \item \textbf{Generator metadata.} Model identifier, prompt
        version, retrieval window, timestamp.
  \item \textbf{Revocation banner} (if revoked). The opinion is
        hidden from the feed but its row and audit trail remain
        readable for accountability.
\end{itemize}

\screenshot{04_currents_audit.png}{The audit trail block at the bottom of a Currents opinion.}{cu-audit}

\section{Follow-up Q\&A on an opinion}

Each opinion has a \page{Follow-up} thread attached. Readers (no
account required) can ask a question or push back on the opinion. The
firm's responses are answered by the Oracle (Guide 3) with the same
citation contract.

This is where most outside engagement happens. Readers who want their
response to be formally considered for the firm's record use the
\page{Public response} flow instead (next section).

\section{Public responses --- the formal reader-engagement channel}

\route{/responses} (public) is the form for a reader to submit a
structured reply: a counter-evidence claim, a counter-argument, a
clarification request, or a request to extend agreement. The form
asks for the article slug, the targeted claim, and the response
body.

A reader who wants their response to be \emph{publishable by name}
ticks a publish-consent box; otherwise the firm holds the response as
internal feedback. The publish-consent is symmetric: even if the firm
later wants to quote the reply publicly, it still needs a fresh
acknowledgment on the founder's reply.

As a founder you'll see the resulting moderation queue on the
operator side (Guide 6); for general use you can read public replies
and engage in follow-up threads.

\section{Critique submissions --- invited expert critique}

\route{/critiques} (public) is a higher-effort form for an invited
external expert to challenge a specific published claim. Submitters
provide their identity (name, email, optional ORCID, optional public
URL), the article slug, the target claim, the counter-evidence, and a
derivation of the alternative position.

Accepted high-severity critiques can earn a small bounty (the firm
pays for being shown wrong, in public). The actual moderation and
bounty confirmation is operator-only --- but the public
\route{/critiques} page lists accepted critiques with submitter credit
(subject to the submitter's hall-of-fame consent).

This is one of the firm's strongest invitations to be wrong: it pays
for being shown wrong, in public.

\section{Subscriber digest}

\route{/subscribe} is the public sign-up. The form is double-opt-in:
the subscriber gets a confirmation email and must click through
before any digest is sent. Subscribers choose their scope
(\texttt{firm}, \texttt{methodology}, \texttt{domain}, or
\texttt{conclusion}) and cadence (\texttt{weekly}, \texttt{immediate},
\texttt{monthly}).

The digest has one operational property worth knowing: \emph{no
tracking pixels}. If a subscriber wants to confirm they read a
digest, they click a one-time link. Otherwise the firm has no idea
who read what --- on purpose.

\section{Articles --- long-form public essays}

The article dispatcher periodically clusters recent opinions and, if
a cluster is dense enough, drafts a long-form public article. Once an
operator has approved the article (the gate is described in Guide 6),
the article lands on \route{/post/<slug>} (when tied to a single
upload) or \route{/c/<slug>/v/<version>} (when promoted from a
reviewed conclusion). A published article carries:

\begin{itemize}[leftmargin=*,itemsep=0.2em]
  \item Its frozen body and citations, so old versions don't change
        when the underlying corpus shifts.
  \item A cryptographic signature attached to the canonical hash,
        written by the workshop and stored on the website. (Verifiers
        can recompute the hash to check the signature.)
  \item An optional \page{Addendum} chain --- dated, never-edits-the-
        original notes that the firm appends quarterly during the
        methodology review week.
\end{itemize}

A reader-side feed (\route{/atom.xml} and \route{/feed.xml}) keeps an
Atom and an RSS view of new articles and reviewed conclusions.

\section{The founder's read of Currents}

As a founder using the platform, the things you typically do with
Currents are:

\begin{itemize}[leftmargin=*,itemsep=0.2em]
  \item Read recent opinions to develop a feel for the firm's voice.
  \item Open the audit trail on any opinion to see which conclusions
        it relied on --- this is useful as a feedback loop on which
        of your conclusions are doing real work.
  \item Watch for abstentions on topics you care about. A recurring
        abstention is a signal that the firm should be reading more
        in that area.
  \item Engage in follow-up threads.
\end{itemize}

\section{Common failure modes (what to expect)}

\begin{itemize}[leftmargin=*,itemsep=0.3em]
  \item \textbf{No opinions in 24 hours.} Either the ingestion path
        is paused or the corpus is thin on relevant topics. Either
        way, founder-alpha sees this on the operations card before
        you do.
  \item \textbf{Opinions look like they are recapping rather than
        opining.} The generator falls back to thinner prose when its
        retrieval comes back near-empty. Solve it by adding firm
        material on the topic.
  \item \textbf{An opinion seems to misread the source.} The
        strawman-detection guard catches most of these but not all.
        Use the \page{Critiques} form to flag it.
\end{itemize}

\section{Where the docs are}

\begin{itemize}[leftmargin=*,itemsep=0.2em]
  \item The other five guides in this series.
  \item Guide 6 covers the operator surfaces for Currents
        (revocation, kill switch, hold for review).
\end{itemize}

\begin{nextup}
Continue with \href{05_Forecasts_and_Portfolio.pdf}{Guide 5 ---
Forecasts and Portfolio} for the prediction-market surface, which
shares Currents' grounding principles but adds a much stricter
operator confirmation flow before any real-money order can be placed.
\end{nextup}

\end{document}
